\begin{circuitikz}[
        american, 
        >={Latex[length=5mm, width=2mm]},
        %>=latex,
        alt={A ladder logic diagram bounded by two vertical power rails, illustrating a seal-in circuit maintaining power after Manual PB is released. The top rung remains unhighlighted and inactive, with Start PB open and output coil LT1 off. On the second rung, the upper branch containing Manual PB is unhighlighted and open. The green highlight now routes from the left rail through the lower parallel branch containing the closed LT2 contact, continuing across the main rung to keep the LT2 output coil energized on the right rail.}
    ]

    \def\rungs{3}   % Power rails are automatically drawn the proper height for # rungs
    \pgfmathsetmacro{\railheight}{-\rungs * 2}

    %========================================================
    % Component Grid
    % Spacing = 3x, 2y
    %========================================================
    \foreach \i in {0,...,5} {
        \foreach \j in {0,...,20} {
        
        % Calculate spatial positions
        \pgfmathsetmacro{\px}{\i * 3}
        \pgfmathsetmacro{\py}{-1 - (\j * 2)}
        
        % Name coordinate using loop indices: (C-column-row) e.g., (C-0-0), (C-5-10)
        \coordinate (C-\i-\j) at (\px, \py);
        
        % OPTIONAL: Uncomment the 2 lines below to display grid labels while building
        % \fill[red] (C-\i-\j) circle (1.5pt);
        %\node[anchor=south west, scale=0.5, gray] at (C-\i-\j) {(\i,\j)};
        }
    }

    % Put the primary logic elements on this layer
    \begin{pgfonlayer}{main}

        \draw[
            BakerBlack
        ] 
        (C-0-0) 
        to[C, color=BakerBlack, font=\small, l={Start PB}] (C-1-0) 
        -- (C-4-0) 
        to[esource, color=BakerBlack, font=\small, l={LT1}] (C-5-0);

        \draw[
            BakerBlack
        ] 
        (C-0-1)
        to[C, color=BakerBlack, font=\small, l={Manual PB}] (C-1-1)
        -- (C-4-1)
        to[esource, color=BakerBlack, font=\small, l={LT2}] (C-5-1);

        \draw[
            BakerBlack
        ]
        (C-0-2)
        to[C, color=BakerBlack, font=\small, l={LT2}] (C-1-2)
        -- (C-1-1);

    \end{pgfonlayer}

    % Use only for green highlighting of true logic elements
    \begin{pgfonlayer}{bg}
        
        % Always highlight the left rail green
        \draw[
            UpNorthGreen,
            opacity=0.6,
            line width=10pt
        ]
        (0,0)
        -- (0,\railheight);

        \draw[
            UpNorthGreen,
            opacity=0.6,
            line width=10pt
        ]
        (C-0-2)
        -- (C-1-2)
        -- (C-1-1)
        -- (C-5-1);

    \end{pgfonlayer}

    % Power rails should be the only thing in the foreground
    % Automatically drawn based on the number of rungs - should need to change this
    % Rails start at y=0 and go down - note the (-) sign is in the macro above
    \begin{pgfonlayer}{ft}
        
        \draw[
            BakerBlack,
            very thick,
        ]
        (0,0)
        -- (0,\railheight);

        \draw[
            BakerBlack,
            very thick,
        ]
        (15,0)
        -- (15,\railheight);

    \end{pgfonlayer}

\end{circuitikz}